\chapter{Rezultate Experimentale}
\label{ch:results}

% One subsection tree per dataset. Each dataset section follows the same
% internal shape for comparability: summary table -> per-model discussion ->
% dataset-specific takeaway. Cross-dataset synthesis (which paradigm wins
% where, and why) is saved for the Discussion/Conclusion chapter, not here --
% keep this chapter to "what happened", not "what it means for the thesis".

\section{Detectarea Fraudei cu Cardul de Credit (Setul de Date A)}
% Documente sursă: data/results/results_credit_card/*_report.txt, credit_card/credit_card_results.md

\subsection{Comparație Generală}

\begin{table}[h]
\centering
\footnotesize
\begin{tabular}{@{}lllll@{}}
\toprule
Model & Paradigmă & ROC-AUC & Precizie Medie & F1 (fraudă) \\
\midrule
MLP Autoencoder            & Reconstrucție              & 0.9649 & \textbf{0.7105} & 0.6992 \\
Isolation Forest           & Partiționare                   & 0.9587 & 0.6725          & 0.6070 \\
LSTM Autoencoder           & Reconstrucție (pseudo-secv.) & 0.9476 & 0.6223          & \textbf{0.7362} \\
OC-SVM                     & Frontieră                     & 0.9354 & 0.5959          & 0.6454 \\
Transformer Bottleneck AE  & Reconstrucție (pseudo-secv.) & 0.9583 & 0.5947          & 0.6137 \\
Predictive LSTM            & Predictiv (pseudo-secv.)     & 0.9581 & 0.5190          & 0.5244 \\
\bottomrule
\end{tabular}
\caption{Detectarea fraudei cu cardul de credit -- comparație de modele, clasificate după Precizia Medie.}
\label{tab:cc-results}
\end{table}

Toate cele șase modele ating ROC-AUC $>0.93$, dar AP variază între 0.519--0.711 $\to$ ROC-AUC este mult mai puțin sensibil la dezechilibrul de clasă față de AP.

\subsection{Discuție}

\paragraph{Nu există o secvență autentică în acest set de date.} Spre deosebire de Yahoo S5 sau variantele temporale/contextuale ale CIC-IDS2017, rândurile de tranzacții nu au nicio ordonare temporală sau cauzală între caracteristici: \texttt{Time} este eliminată la preprocesare (\S\ref{sec:preprocessing}); modelele secvențiale au fost adăugate ulterior exclusiv pentru comparabilitate între seturile de date.
Cele trei arhitecturi \enquote{temporale} evaluate aici (Predictive LSTM, LSTM AE, Transformer Bottleneck AE) nu văd o secvență reală -- ele reformatează fiecare rând cu 29 de caracteristici într-o secvență artificială de 29 de pași scalari ($(N,29) \to (N,29,1)$), câte o caracteristică pe pas, în ordinea din fișierul .csv.

\paragraph{Impact asupra rezultatelor.} Predictive LSTM este cel mai slab model (AP 0.519): obiectivul său de predicție a următorului pas presupune o relație cauzală între pași consecutivi, dar componentele PCA adiacente sunt ortogonale/decorelate prin construcție, deci nu există niciun semnal real de învățat -- reflectat în cel mai slab număr de fals-pozitive (FP=566).
LSTM AE se descurcă comparativ bine (AP 0.622, cel mai bun F1 dintre toate modelele, 0.736, cel mai mic FP, 174), probabil întrucât encoderul său comprimă întreaga intrare într-o singură stare ascunsă sumară, funcționând similar bottleneck-ului MLP AE, mai degrabă decât depinzând doar de ordine.
Transformer Bottleneck AE (AP 0.595) se situează aproape de modelele clasice -- apare agregarea medie globală (global average pooling) într-un vector latent.
MLP Autoencoder, singura arhitectură care tratează cele 29 de caracteristici ca vector plat obișnuit, câștigă per total (AP 0.7105) -- consistent cu absența oricărei secvențe de exploatat. Impunerea unei structuri secvențiale acolo unde nu există niciuna pare să fie un handicap.

\paragraph{Observații la nivel de paradigmă.} Reconstrucția este cea mai puternică paradigmă per total, odată ce se ia în calcul handicapul pseudo-secvențial (MLP AE, LSTM AE și Transformer Bottleneck AE ocupă trei din primele patru locuri).
Isolation Forest este cel mai bun model din afara paradigmei de reconstrucție (AP 0.6725) și al doilea per total. OC-SVM (AP 0.5959) se situează la mijloc, cu al doilea cel mai mare număr de fals-pozitive după Predictive LSTM, singurul reprezentant al paradigmei predictive aici și cel mai slab model per total, din motivul structural discutat mai sus.

\paragraph{Sanity check.} F1-ul cel mai bun dintre toate (0.7362) al LSTM AE, în ciuda unui AP mediocru (0.6223), reamintește că F1 este calculat la un singur prag oracol F2 și poate favoriza un punct de operare diferit de cel care maximizează aria de sub curba PR.



\section{Yahoo Webscope S5 (Setul de Date B)}
% Documente sursă: yahoo/results_notes.md, data/results/results_yahoo/*_report.txt

A2 este exclus din comparațiile de mai jos: la fel ca A1, conține doar anomalii punctuale, fără a adăuga un tip nou de anomalie în comparație.

\subsection{A1 -- Anomalii Punctuale (Metrici Reale de Server)}

Causal Transformer și Transformer Bottleneck AE nu au fost rulate pe acest benchmark.

\begin{table}[h]
\centering
\footnotesize
\begin{tabular}{@{}llll@{}}
\toprule
Model & Paradigmă & ROC-AUC & Precizie Medie \\
\midrule
MLP Autoencoder  & Reconstrucție & 0.792 & \textbf{0.611} \\
LSTM Autoencoder & Reconstrucție & 0.794 & 0.600 \\
OC-SVM           & Frontieră       & 0.776 & 0.595 \\
Predictive LSTM  & Predictiv     & 0.789 & 0.588 \\
Isolation Forest & Partiționare      & 0.765 & 0.511 \\
\bottomrule
\end{tabular}
\caption{Yahoo S5 A1 (anomalii punctuale) -- comparație de modele, clasificate după Precizia Medie.}
\label{tab:yahoo-a1}
\end{table}

\paragraph{Discuție.} Toate cele cinci modele se grupează strâns în intervalul AP 0.51--0.61: MLP Autoencoder este cel mai bun (AP 0.611), Isolation Forest cel mai slab (AP 0.511), dar diferența e suficient de mică încât nicio arhitectură nu arată un avantaj structural aici -- anomaliile punctuale (vârfuri pe un singur pas temporal) nu necesită înțelegere temporală pentru a fi detectate.

\subsection{A3 -- Anomalii Contextuale/Sezoniere}

\begin{table}[h]
\centering
\footnotesize
\begin{tabular}{@{}llll@{}}
\toprule
Model & Paradigmă & ROC-AUC & Precizie Medie \\
\midrule
Predictive LSTM            & Predictiv     & \textbf{0.868} & \textbf{0.812} \\
Transformer Bottleneck AE  & Reconstrucție & 0.762          & 0.663 \\
Causal Transformer         & Predictiv     & 0.738          & 0.623 \\
MLP Autoencoder            & Reconstrucție & 0.765          & 0.507 \\
LSTM Autoencoder           & Reconstrucție & 0.566          & 0.392 \\
OC-SVM                     & Frontieră       & 0.584          & 0.376 \\
Isolation Forest           & Partiționare      & 0.511          & 0.331 \\
\bottomrule
\end{tabular}
\caption{Yahoo S5 A3 (anomalii contextuale/sezoniere) -- comparație de modele, clasificate după Precizia Medie.}
\label{tab:yahoo-a3}
\end{table}

\paragraph{Discuție.} Rezultatele diverg puternic odată ce anomaliile devin contextuale -- o valoare normală în termeni absoluți, dar greșită pentru faza sa în ciclul sezonier.
Isolation Forest și OC-SVM se prăbușesc aproape de random (ROC-AUC 0.511 și, respectiv, 0.584).
LSTM Autoencoder se prăbușește la fel de rău (ROC-AUC 0.566), în ciuda faptului că este un model secvențial -- ceea ce arată că este vorba de o eșuare cauzată de blocajul de compresie (bottleneck), nu de o slăbiciune a arhitecturilor recurente în general.
În aceeași ordine de idei, Transformer Bottleneck AE rămâne în urma celui mai bun model (AP 0.663), dar este de departe cea mai puternică arhitectură bazată pe \emph{reconstrucție}, cu mult înaintea LSTM AE: auto-atenția globală asupra ferestrei recuperează probabil periodicitatea pe care blocajul LSTM AE o distruge, ceea ce probabil explică diferența.
Predictive LSTM domină și este cel mai bun model pe orice benchmark Yahoo (AP 0.812, ROC-AUC 0.868): obiectivul său de predicție a pasului următor exploatează direct predictibilitatea unei serii sezoniere, astfel încât o anomalie injectată produce un vârf clar în eroarea de predicție exact acolo unde violează traiectoria așteptată.

\subsection{A4 -- Anomalii de Tip Schimbare-de-Nivel}

\begin{table}[h]
\centering
\footnotesize
\begin{tabular}{@{}llll@{}}
\toprule
Model & Paradigmă & ROC-AUC & Precizie Medie \\
\midrule
Predictive LSTM            & Predictiv     & \textbf{0.745} & \textbf{0.505} \\
Causal Transformer         & Predictiv     & 0.695          & 0.436 \\
MLP Autoencoder            & Reconstrucție & 0.719          & 0.402 \\
Transformer Bottleneck AE  & Reconstrucție & 0.655          & 0.404 \\
Isolation Forest           & Partiționare      & 0.498          & 0.268 \\
OC-SVM                     & Frontieră       & 0.519          & 0.265 \\
LSTM Autoencoder           & Reconstrucție & 0.502          & 0.252 \\
\bottomrule
\end{tabular}
\caption{Yahoo S5 A4 (anomalii de tip schimbare-de-nivel) -- comparație de modele, clasificate după Precizia Medie.}
\label{tab:yahoo-a4}
\end{table}

\paragraph{Discuție.} Seriile modelează acum un sistem care experimentează schimbări structurale bruște, iar fiecare model se degradează față de A3, indiferent de paradigmă. Un aspect important: ferestrele din toate cele $\sim$100 de serii sunt grupate laolaltă fără identitatea seriei, iar fiecare serie se situează la propriul nivel de bază, deci valoarea absolută a unei ferestre, luată singură, nu poate distinge o fereastră autentică de după schimbare de o fereastră obișnuită dintr-o altă serie al cărei nivel de bază e similar -- de unde rezultatele mai slabe ale modelelor bazate pe distribuția agregată.
Modelele clasice pierd mult sub această confuzie -- Isolation Forest ajunge la ROC-AUC 0.498. Autoencoderele de reconstrucție (MLP AE, Transformer Bottleneck AE) sunt afectate din același motiv și nu depășesc Predictive LSTM aici, în ciuda faptului că o fereastră în formă de treaptă ar fi intuitiv mai ușor de semnalat prin eroarea de reconstrucție.
Predictive LSTM încă conduce (AP 0.505, ROC-AUC 0.745), deși marja sa se îngustează puternic față de A3.

\subsection{Verificare de Referință Trivială (Doar Yahoo)}
\label{sec:yahoo-trivial-baseline}

Urmând \textcite{kim2022rigorous}, care recomandă verificarea fiecărui model antrenat față de cel puțin o referință fără etichete (label-free), modelele Yahoo S5 au fost evaluate suplimentar față de o astfel de referință: norma L2 brută a ferestrei $\lVert w \rVert_2$, fără niciun model învățat. Toate modelele antrenate depășesc clar această referință pe A3, confirmând că avantajele observate mai sus sunt reale, nu un artefact al magnitudinii ferestrei. Scorurile referinței și marja exactă per model sunt prezentate în Anexa~\ref{sec:yahoo-trivial-baseline-appendix}.

\subsection{Discuție}

\paragraph{Atât paradigma, cât și arhitectura contează.} Rezultatele A3 formează o comparație de tip 2$\times$2 între paradigmă (reconstrucție vs. predictiv) și arhitectură (LSTM vs. Transformer): reconstrucția obține 0.39 AP (LSTM AE) și 0.66 AP (Transformer Bottleneck), paradigma predictivă obține 0.81 AP (Predictive LSTM) și 0.62 AP (Causal Transformer).
Paradigma predictivă este necesară -- ambele modele predictive depășesc ambele modele de reconstrucție -- dar nu suficientă de una singură: diferența de 0.19 AP dintre Predictive LSTM și Causal Transformer, cu un obiectiv de antrenare identic, arată că arhitectura LSTM poartă un avantaj suplimentar pe date sezoniere. Pe A4, acest decalaj se reduce la 0.07 AP -- o singură tranziție bruscă nu necesită acumulare de fază, deci biasul inductiv secvențial al LSTM-ului contează mai puțin.

\paragraph{Rezumat pe tip de anomalie.} Anomaliile punctuale (A1) lasă toate modelele într-o bandă îngustă de AP 0.51--0.61 -- nu e necesară nicio înțelegere temporală pentru a semnala un vârf izolat.
Anomaliile sezoniere (A3) sunt locul unde avantajul paradigmei predictive e cel mai clar, și unde alegerea arhitecturii contează cel mai mult în cadrul acelei paradigme. Punctele de schimbare (A4) degradează fiecare model, dar Predictive LSTM încă conduce, iar Causal Transformer încă depășește ambele modele de reconstrucție. Modelele clasice (Isolation Forest, OC-SVM) și LSTM Autoencoder se prăbușesc pe orice tip de anomalie care necesită discriminare contextuală, nu absolută (A3, A4), rămânând competitive doar pe A1.



\section{CIC-IDS2017 (Setul de Date C)}
% Documente sursă: cicids2017/results_analysis.md

\subsection{Comparație Generală}

Fiecare valoare AP raportată mai jos este calculată pe întreaga mulțime de test flat/context agregată, în care PortScan, DoS Hulk și DDoS reprezintă împreună 95.1\% din toate fluxurile de atac (compoziția completă pe clase: Tabelul~\ref{tab:cic-composition}, Anexa~\ref{ch:cic-composition-temporal}) -- susținând afirmația, făcută repetat mai jos, că AP-ul agregat e practic un proxy pentru performanța pe aceste trei clase. Celelalte zece clase (toate sub 2\% individual) au o influență limitată asupra clasamentului modelelor, de unde necesitatea unei detalieri pe atac (secțiunea~\ref{sec:cic-per-attack}).
\textbf{DoS Hulk} (158{,}469 fluxuri) rămâne parte a clasei pozitive pentru fiecare metrică agregată, dar e exclus din detalierea \emph{pe clasă} din secțiunea~\ref{sec:cic-per-attack}, întrucât \textcite{engelen2021troubleshooting} arată că unealta folosită pentru a-l genera e configurată greșit (setează antetul HTTP \texttt{Connection} pe \texttt{Close} în loc de \texttt{Keep-Alive}), făcând atacul ineficient prin construcție.

\begin{table}[H]
\centering
\resizebox{0.72\textwidth}{!}{%
\begin{tabular}{@{}lllllll@{}}
\toprule
Model & Set de caracteristici & Paradigmă & ROC-AUC & Precizie Medie & Acuratețe Echilibrată & F1 (atac) \\
\midrule
MLP Autoencoder            & Context  & Reconstrucție & \textbf{0.9860} & \textbf{0.9487} & 0.9548          & 0.9020 \\
OC-SVM                     & Context  & Frontieră       & 0.9745          & 0.9396          & \textbf{0.9569} & \textbf{0.9229} \\
Isolation Forest           & Context  & Partiționare      & 0.9785          & 0.9328          & 0.9270          & 0.8367 \\
LSTM Autoencoder           & Temporal & Reconstrucție & 0.8976          & 0.7955          & 0.7846          & 0.6249 \\
MLP Autoencoder            & Flat     & Reconstrucție & 0.9189          & 0.7921          & 0.8022          & 0.6305 \\
Isolation Forest           & Flat     & Partiționare      & 0.9049          & 0.7612          & 0.8228          & 0.6639 \\
OC-SVM                     & Temporal & Frontieră       & 0.8629          & 0.7816          & 0.7600          & 0.5888 \\
MLP Autoencoder            & Temporal & Reconstrucție & 0.8831          & 0.7509          & 0.7734          & 0.6125 \\
Isolation Forest           & Temporal & Partiționare      & 0.8438          & 0.7604          & 0.6841          & 0.5195 \\
OC-SVM                     & Flat     & Frontieră       & 0.7336          & 0.6845          & 0.6121          & 0.4644 \\
Bottleneck Transformer AE  & Temporal & Reconstrucție & 0.8077          & 0.6539          & 0.7059          & 0.5423 \\
Causal Transformer         & Temporal & Predictiv     & 0.7099          & 0.5449          & 0.5520          & 0.4320 \\
Predictive LSTM            & Temporal & Predictiv     & 0.6126          & 0.4472          & 0.5475          & 0.4295 \\
\bottomrule
\end{tabular}%
}
\caption{CIC-IDS2017 -- toate combinațiile model/set-de-caracteristici, clasificate după Precizia Medie.}
\label{tab:cic-results}
\end{table}

\paragraph{Caracteristicile de context al gazdei domină.} Cele trei modele augmentate cu context ocupă primele trei poziții la fiecare metrică, cu AP între 0.933--0.949. Diferența dintre cel mai bun model cu context (AE+Context, AP 0.949) și cel mai bun model fără context (MLP AE Flat, AP 0.792) depășește întreaga plajă acoperită de celelalte zece modele.
Fiecare dintre cele trei arhitecturi de bază se îmbunătățește substanțial la adăugarea caracteristicilor de context: IForest câștigă +0.172 AP (0.761 $\to$ 0.933), MLP AE câștigă +0.157 (0.792 $\to$ 0.949), iar OC-SVM câștigă cel mai mult, +0.255 (0.685 $\to$ 0.940) -- semn că uneori caracteristicile mai bune contează mai mult decât alegerea paradigmei.

\paragraph{Adăugarea ferestrelor temporale, singură, nu ajută.} Pentru IForest și MLP AE, gruparea fluxurilor în ferestre ordonate după IP-ul sursă (secțiunea~\ref{sec:preprocessing}) nu oferă o îmbunătățire semnificativă față de vederea flat, și uneori chiar dăunează -- contrast direct cu câștigul din contextul gazdei: ambele modificări adaugă informație despre \emph{ordinea/activitatea fluxurilor în timp}, dar doar statisticile de context agregate se traduc într-un semnal utilizabil pentru cele două clase de atac (DDoS, PortScan) care domină setul de test ca volum.

\paragraph{Reconstrucția depășește predicția.} Printre modelele cu caracteristici temporale, arhitecturile bazate pe reconstrucție (LSTM AE, AE Temporal, OC-SVM Temporal) depășesc clar ambele modele bazate pe predicție (Predictive LSTM AP 0.447, Causal Transformer AP 0.545) -- inversul constatării de la Yahoo A3 (capitolul~\ref{ch:results}). Sweep-urile de prag ale ambelor modele predictive ajung la un punct de operare extrem (FP $>$ 1.1 mil. din 1.28 mil. de fluxuri benigne), consistent cu ROC-AUC-ul lor scăzut (0.61 și 0.71).

\subsection{Detaliere pe Tip de Atac}
\label{sec:cic-per-attack}

\begin{table}[h]
\centering
\resizebox{0.78\textwidth}{!}{%
\begin{tabular}{@{}lrlll@{}}
\toprule
Clasă de atac & $n$ & IForest+Ctx & OC-SVM+Ctx & AE+Ctx \\
\midrule
DDoS                & 95{,}000  & 0.982          & \textbf{0.998} & 0.994          \\
PortScan             & 159{,}000 & 0.601          & 0.820          & \textbf{0.842} \\
DoS GoldenEye        & 7{,}500   & \textbf{0.655} & 0.215          & 0.220          \\
DoS slowloris        & 4{,}000   & \textbf{0.190} & 0.012          & 0.025          \\
DoS Slowhttptest     & 1{,}700   & \textbf{0.112} & 0.007          & 0.013          \\
Heartbleed           & 11        & \textbf{0.093} & 0.001          & 0.000          \\
Infiltration         & 32        & \multicolumn{3}{l}{AP 0.03--0.14, fără un câștigător clar la niciun model} \\
Bot                  & 738       & \multicolumn{3}{l}{$<$0.01 la toate modelele, inclusiv variantele cu context} \\
FTP-Patator          & 4{,}000   & \multicolumn{3}{l}{$<$0.01 la toate modelele, inclusiv variantele cu context} \\
SSH-Patator          & 3{,}000   & \multicolumn{3}{l}{$<$0.01 la toate modelele, inclusiv variantele cu context} \\
Subtipuri Web Attack & 12--151   & \multicolumn{3}{l}{AP aproape zero peste tot (Brute Force, SQL Injection, XSS)} \\
\bottomrule
\end{tabular}%
}
\caption{CIC-IDS2017 -- Precizia Medie pe clasă de atac pentru cele trei modele augmentate cu context.}
\label{tab:cic-per-class}
\end{table}

\paragraph{Detectate constant.} Două clase de atac ies în evidență ca fiind detectabile fiabil pe aproape orice set de caracteristici:

\begin{itemize}
  \item \textbf{DDoS} este cea mai ușoară clasă, cu o marjă mare: toate modelele cu context obțin AP $>$ 0.98 (OC-SVM+Context atinge 0.998), iar chiar și modelele flat depășesc AP 0.77 -- posibil pentru că rata mare de pachete mici, tipică unei inundații, produce o valoare aberantă extremă în ambele spații de caracteristici.
  \item \textbf{PortScan} este cea mai clară demonstrație a valorii caracteristicilor de context: fiecare model flat și temporal obține AP 0.07--0.28, în timp ce fiecare model cu context depășește 0.60, AE+Context atingând 0.842. Un singur flux de scanare scurt, către un port arbitrar, e nesemnificativ izolat, dar numărul de porturi de destinație unice contactate de IP-ul sursă în ultimele 60 de secunde ajută probabil în special la acest atac. IForest+Context e comparativ mai slab aici (0.601).
\end{itemize}

\paragraph{Răspuns mixt la caracteristici de context pentru DoS la nivel de aplicație.} Nu fiecare clasă de atac răspunde la caracteristicile de context în același mod -- cele trei clase DoS la nivel de aplicație împart familia de context în loc să o favorizeze uniform:
\textbf{DoS GoldenEye} arată o separare clară: IForest+Context atinge AP 0.655, iar AE+Context și OC-SVM+Context se prăbușesc amândouă la $\sim$0.22.
\textbf{DoS Slowhttptest și DoS slowloris} arată același tipar, mai atenuat: IForest+Context e singurul model cu context cu detecție semnificativă (AP 0.112 și 0.190), rămânând totuși mult sub AP-ul de 0.69/0.61 al IForest Temporal pe aceleași clase, probabil întrucât fereastra de context de 60s e o nepotrivire parțială pentru conexiunile lente, de minute întregi, ale acestor atacuri -- ordonarea cronologică explicită per-gazdă a IForest Temporal rămâne cea mai bună detecție pentru această familie.

\paragraph{În mare parte nedetectabile, indiferent de setul de caracteristici.} FTP-Patator, SSH-Patator, Bot și cele trei subtipuri Web Attack obțin toate AP $<$ 0.01 la fiecare model, inclusiv cele cu context. Detectarea lor ar necesita cel mai probabil contoare de autentificare eșuată sau inspecția payload-ului, niciuna disponibilă în acest set. Heartbleed (11 eșantioane) și Infiltration (32) sunt prea mici pentru concluzii solide.

\subsection{Discuție}

\paragraph{Caracteristicile de context al gazdei sunt cea mai impactantă adăugare din acest studiu.} Adăugarea a 9 statistici de activitate per-gazdă cu fereastră glisantă (secțiunea~\ref{sec:preprocessing}) ridică AP-ul de la 0.79 la 0.95 per ansamblu -- un câștig mai mare decât alegerea paradigmei sau trecerea de la o vedere flat la una temporală -- determinat în principal de rezolvarea punctului orb PortScan comun tuturor modelelor flat și temporale. Niciun model de context, luat singur, nu domină la fiecare tip de atac, deși AE+Context câștigă la AP-ul agregat, determinat de cele două clase cele mai mari.

\paragraph{Un ansamblu practic.} Niciun model sau set de caracteristici, luat singur, nu acoperă fiecare clasă de atac din acest set de date. Combinarea AE+Context (puternic pe DDoS și PortScan, cele două clase care domină setul de test ca volum) cu IForest Temporal (puternic pe clasele DoS lente și, cu titlu provizoriu, pe Heartbleed) ar acoperi cea mai largă gamă de clase de atac detectabile în cadrul acestui set de caracteristici. Această premisă, precum și un ansamblu cu scor combinat construit și evaluat pe baza ei, sunt testate direct în secțiunea~\ref{sec:explain-disjoint-regime}: evidența celor două modele este confirmată ca fiind cu adevărat disjunctă, dar combinarea scorurilor lor se dovedește a schimba calitatea detecției PortScan pentru câștiguri pe clasele DoS lente, mai degrabă decât să le acopere pe amândouă gratuit.